%\documentclass[notes=show]{beamer}
\documentclass[notes=none]{beamer}

\usepackage{tikz}
\usepackage[linguistics]{forest}
\usepackage[many]{tcolorbox}
\usepackage{comment}
\usepackage{multicol}
\usepackage{listings}

\newif\ifframeinlbf
\frameinlbftrue
\makeatletter
\newcommand\listofframes{\@starttoc{lbf}}
\makeatother

\addtobeamertemplate{frametitle}{}{%
    \ifframeinlbf%
         \addtocontents{lbf}{%
            \protect\makebox[2em][l]{%
                \protect\usebeamercolor[fg]{structure}\insertframenumber\hfill%
            }%
            \protect\hyperlink{page.\insertpagenumber}\insertframetitle\par%
        }%
   \else\fi%
}

\AtBeginSection[]{
  \begin{frame}
  \vfill
  \centering
  \begin{beamercolorbox}[sep=8pt,center,shadow=true,rounded=true]{title}
    \usebeamerfont{title}\insertsectionhead\par%
  \end{beamercolorbox}
  \vfill
  \end{frame}
}

\newcommand{\myalert}[2]{%
    \tikz\node[coordinate](t#1){};\alert<#1>{#2}
}


\newcommand{\myexplain}[3]
{\begin{tcolorbox}
    [   tikznode boxed title,
        enhanced,
        arc=0mm,
        interior style={white},
        attach boxed title to top center= {yshift=-\tcboxedtitleheight/2},
        fonttitle=\bfseries,
        colbacktitle=white,coltitle=red,
        boxed title style={size=small,colframe=red,boxrule=0pt},
        title={\tikz\node[coordinate](s#1) {};#2}]
#3
\end{tcolorbox}
\begin{tikzpicture}[overlay]
\path<#1>[->, thick] (s#1) edge [bend left] (t#1);
\end{tikzpicture}
}

\newcommand{\explain}[2]
{\begin{tcolorbox}
    [   tikznode boxed title,
        enhanced,
        arc=0mm,
        interior style={white},
        attach boxed title to top center= {yshift=-\tcboxedtitleheight/2},
        fonttitle=\bfseries,
        colbacktitle=white,coltitle=red,
        boxed title style={size=small,colframe=red,boxrule=0pt},
        title={#1}]
#2
\end{tcolorbox}
}

\def\point{\begin{tikzpicture}
\fill[orange] (0,0) circle (2pt);
\end{tikzpicture}}

\setbeamertemplate{navigation symbols}{} %Remove navigation bar

\tikzstyle{every picture}+=[remember picture]

\usepackage{color}
\newcommand{\vet}[1]{\foreach \num in {#1}{\el{\num}}}
\newcommand{\el}[1]{\fbox{\parbox[c][2ex][c]{1.1em}{#1}}}
\newcommand{\id}[1]{\parbox[r][2ex][b]{1.7em}{\raggedleft #1}}
\newcommand{\bl}{\color{blue}}
\newcommand{\gr}{\color{gray}}
\newcommand{\re}{\color{red}}

\title{Languages}
\author{Lunjin Lu}
\date{}


\begin{document}
\begin{comment}
\end{comment}
\frameinlbffalse
\begin{frame}[fragile,t,allowframebreaks]{List of Frames}
    \listofframes
\end{frame}
\frameinlbftrue


\frame{ \titlepage
}

\begin{frame}{Languages}
\begin{itemize}
  \item There are numerous languages - natural languages, programming languages, ...
  \item Natural languages - English, French, Chinese, Japanese, ...
  \item Programming languages - Java, C++, Prolog, F\#, ...
\end{itemize}
\end{frame}

\begin{frame}{Language Examples}
    \begin{itemize}
    \item The English language is the set of all English sentences composed of English words and punctuation marks.  Its alphabet is the set of English words and punctuation marks.
    \item The Java language is the set of all valid Java programs. Its alphabet contains keywords, operators, identifiers, integers, ...
    \end{itemize}
    Members of the alphabet of a language are  called called \textit{symbols}; and members of the language are called \textit{strings}.
\end{frame}

\begin{frame}{Strings}
A string over an alphabet $\Sigma$ is a sequence of symbols from $\Sigma$. The empty string is denoted $\Lambda$.



Let $\Sigma = \{0,1\}$. The following are example strings over $\Sigma$ include
\begin{itemize}
    \item $\Lambda$
    \item 0
    \item 1
    \item 100
    \item 10000
\end{itemize}

The length of a string $w$, denoted $|w|$ is defined as the number of symbols occurrences in $w$. For instance, $|\Lambda|=0$ and $|010|=3$.
\end{frame}


\begin{frame}{What is a Language?}
    A language $L$  is a set of strings over a given alphabet $\Sigma$.
    \begin{itemize}
    \item [Ex 1.] $\Sigma_1=\{0,1\}$; 
    \[L_1=\left\{\begin{array}{c}
    \mbox{All strings of 0's and 1's such that}\\
    \mbox{if a string starts with 0}\\
    \mbox{then its length is 1} \end{array}\right\}\] 
    \item [Ex 2.] $\Sigma_2=\{a\}$; $L_2=\{\mbox{all strings of a's whose lengths are odd}\}$
    \end{itemize}
\end{frame}

\begin{frame}{String Reversal}
    The reversal of a string $w$, denoted $w^{-1}$,  is the string resulting from the arrangement of  the symbols in $w$ in reverse order. If $w=a_{1}a_{2}...a_{n-1}a_{n}$, then its reversal $w^{-1}=a_{n}a_{n-1}...a_{2}a_{1}$.
    \begin{itemize}
    \item $\Lambda^{-1}=\Lambda$
    \item $0^{-1}=0$
    \item $(abc)^{-1}=cba$
\end{itemize}
\end{frame}

\begin{frame}{String Concatenation}
    The concatenation of two strings $x$ and $y$, written as  $x\cdot y$ (or simply $xy$) is a string formed of symbols from $x$ followed by symbols from $y$ in order. If $x=a_{1}a_{2}...a_{m}$ and $y=b_1b_2\cdots b_n$ then $x\cdot y=a_1a_2...a_mb_1b_2...b_n$.
    \begin{itemize}
    \item $\Lambda\cdot x=x$ and $x\cdot\Lambda=x$ for any string $x$.
    \item $|x\cdot y| = |x|+|y|$.
    \item $101\cdot 001=101001$.
    \item $001\cdot 101=001101$.
\end{itemize}
$\Lambda$ is the identity for concatenation and concatenation is not commutative.
\end{frame}

\begin{frame}{Concatenation of two languages}
Let $L_1$ and $L_2$ be two languages. The concatenation of $L_1$ and $L_2$ is defined
\[ L_1L_2 = \{w_1\cdot w_2 \mid w_1\in L_1\wedge w_2\in L_2\}
\]

Let $L_1=\{a,ab,abb\}$ and $L_2=\{c,cd\}$. Then
\[ L_1L_2 = \{ac,abc,abbc,acd,abcd,abbcd\}
\]
\end{frame}

\begin{frame}{Kleene Closure of a Language}

The Kleene closure of a language $L$, denoted $L^*$, is the language whose elements are concatenations of strings from $L$.

\[ L^* = \{w_1\cdots w_n \mid n\geq 0 \wedge \forall 1\leq i\leq n. w_i\in L\}
\]


Let $L_1=\{1,10\}$. Then $L_1^{*} = \{\Lambda, 1, 10, 11, 110,1010, \cdots \}$.

Let $L_2=\{0,1\}$. Then $L_2^*=\{\mbox{Strings of 0's and 1's}\}$.

Note that $\Lambda \in L^*$ for any language $L$ and $L^*$ is infinite if $L\neq \emptyset$ and $L\neq \{\Lambda\}$.

\end{frame}

\begin{frame}{Proof by Contradiction} 
Proof by contradiction is one of common proof methods. It works as follows.
 
In order to prove a property P, 
\begin{itemize} 
\item first assume that P does not hold, and then, 
\item show that this assumption leads to contradiction.
\end{itemize}
\end{frame}

\begin{frame}{Prove that $\sqrt{2}$ is not a rational number}
Proof. Assume that $\sqrt{2}$ is a rational number. There are two integers 
 	   $p$ and $q$ such that $\sqrt{2}=p/q$ and $p$ and $q$ do not have 
 	   common divisors other than 1. There are prime number $p_1,p_2\cdots p_m$ such that $p=p_1p_2\cdots p_m$ and 
 	   $p_1\leq p_2\leq\cdots\leq p_m$,  and    prime number $q_1,q_2\cdots q_n$ such that $q=q_1q_2\cdots q_n$ and 
 	   $q_1\leq q_2\leq\cdots\leq q_n$. 
 	   
 	   We have $\sqrt{2} = \frac{p_1p_2\cdots p_m}{q_1q_2\cdots q_n}$ which is equivalent to 
 	   $\sqrt{2}{q_1q_2\cdots q_n} = {p_1p_2\cdots p_m}$ 
 	   Multiply each side by itself, we obtain 
 	   $2{q_1^2q_2^2\cdots q_n^2}={p_1^2p_2^2\cdots p_m^2}$. Both sides are products of prime numbers in ascending order. 
 	   Therefore, $p_1=2$ and hence 
 	   ${q_1^2q_2^2\cdots q_n^2}=2p_2^2\cdots p_m^2$. Both sides are again products of prime numbers in ascending order 
 	   and hence $q_1=2$. Thus 2 is a common divisor of $p$ and $q$, which contradicts with that $p$ and $q$ do not have 
 	   common divisor other than 1. Therefore, the assumption is incorrect and hence $\sqrt{2}$ is not a rational number.
\end{frame}

\end{document}
